\section{基于名称相似度与外部语义向量的特征抽取与注入}
\label{sec:sem_encoder_inject}

覆盖与结构证据在许多缺陷版本中能够给出清晰的可疑度差异，但当失败测试触达高度重叠且候选方法处于相近结构邻域时，排序往往出现打平现象。此时最棘手的问题并不是缺少证据，而是不同候选在同一类证据上呈现出近似的形态，导致排序的间隔难以被稳定拉开。本文因此在现有实现中补充一条成本更低但更稳定的语义线索，其来源并非在线调用预训练编码器对方法体进行编码，而是直接从方法名与测试名中提取名称层面的相似度信号，并在条件允许时用外部 compact 语义向量增强这一信号。为了避免语义模块对主干传播逻辑造成侵入式改动，语义证据最终被压缩为一个标量分数，作为节点初始表示中的一维附加特征进入图推断过程，而不是在传播阶段被单独加权或单独建模。

从实现流程看，这条语义线索像一条短路径。它从原始标识符出发，先把难以直接比较的名称字符串转换为可比的 token 表示，再把 token 级相似度压缩为单一的语义分数，最后把这一分数落位到方法与测试节点上，使其与覆盖与结构证据在同一节点表示空间内共同参与后续的结构传播与排序学习。下文按这一过程进行叙述。

\subsection{从标识符到可比 token 的转换}

语义证据的起点是方法与测试的名称信息，但名称字符串本身并不可直接比较。一方面，不同项目的包结构深度与命名习惯差异较大，完整限定名会引入大量与功能无关的路径噪声。另一方面，方法与测试常见驼峰命名、下划线命名或二者混合，且经常伴随缩写、数字后缀与约定俗成的前后缀，若直接以原始字符串做匹配，相似度会被分隔符选择与命名风格强烈左右，从而使语义证据更多反映书写习惯而非语义关联。基于这一考虑，本文在计算名称相似度之前，先将方法与测试的标识符统一转换为可比的 token 序列，使后续相似度计算可以落在一致口径的离散集合运算上。

实现中，方法标识符以 \texttt{Class:Method} 形式给出。对于 \texttt{Class} 部分，本文并不保留完整包路径，而是截取末两级包名作为上下文锚点，并与类名、方法名共同构成用于分词的输入字符串。这样做的直觉是保留少量语境信息以区分同名方法，同时避免深层包路径带来的域特异噪声与跨项目不可比性。随后对 ``.'' 等路径分隔符进行去除或边界化处理，保证包路径不会以原样进入 token 序列，从而避免低信息 token 大量堆积并稀释真正具有区分度的名称片段。

在分词策略上，本文采用与代码一致的 \texttt{splitCamel} 规则对标识符进行归一化分词。该规则同时覆盖驼峰与下划线两类常见命名形式。对于驼峰命名，按大小写边界切分为子词，例如 \texttt{getUserName} 被拆解为 \texttt{get}、\texttt{user}、\texttt{name}，从而把一个复合标识符还原为更接近自然语言的词项集合。对于下划线命名，将 ``\_'' 作为显式边界进行切分，并对切分后的片段继续应用驼峰拆解，从而使 \texttt{get\_user\_name} 与 \texttt{getUserName} 这类同义命名在 token 层面趋于一致。对于包含数字与缩写的标识符，分词过程保留其作为独立片段参与后续重叠统计，以避免把工程中常见的序号、版本后缀或缩写模式完全抹去。测试标识符采用与方法完全相同的归一化与分词规则，使方法 token 与测试 token 在粒度与生成口径上严格一致，从而保证后续相似度值可解释且可比较。通过这一转换，名称相似度被规约为 token 集合之间的重叠统计问题，为语义分数生成提供了稳定输入。

\subsection{从名称相似度到语义分数}

完成 token 化后，语义证据以一个标量分数的形式产生。该设计背后有一个明确取舍。本文并不试图在语义层面替代覆盖与结构证据，而是将语义视作一条轻量补充线索，其价值在于当候选在覆盖与结构维度上趋同打平时，名称层面的关联能够提供额外的微弱偏好，使排序更容易形成稳定间隔。因此，语义证据被压缩为单一数值，既便于与节点嵌入直接融合，又避免在模型内部引入额外的序列编码与对齐开销。

默认设置下，本文在 token 层面计算名称重叠比例，并将其解释为方法与测试在命名层面的语义关联强度。实现中使用 \texttt{getoverlap} 在方法 token 集合与测试 token 集合之间计算重叠比例，并对同一缺陷版本内的所有测试逐一评估，得到一组方法到测试的相似度值。随后取其中的最大值作为该方法的语义分数。最大值聚合体现了故障定位中的局部指示特征。实践中往往只有少量测试对缺陷位置具有强指示性，若对所有测试做均值聚合，强关联会被大量弱关联稀释，语义分数将更像一种全局背景而非可区分线索。最大值聚合则保留了最强关联信号，使语义分数能够更直接地体现某个候选方法与某个测试之间的显著命名联系，从而在候选打平时更可能提供可用的排序偏置。

得到最大重叠比例后，本文将其映射为整数区间 $1\!\sim\!101$，形成 1 维语义标量特征 \texttt{inputtext}。这一离散化处理并非为了提升表达能力，而是为了让语义分数以更稳健的尺度参与融合。一方面，离散化削弱了极小差异对后续学习过程的敏感性，使语义信号更像一个粗粒度提示而非精细回归量。另一方面，固定区间便于与离散嵌入的数值范围并置，降低直接拼接时的尺度失衡风险。由于语义分数只承担补充区分线索的角色，其关键在于版本内的相对可比性而非连续精度，因此离散映射在工程上更稳健。

在默认名称重叠之外，框架支持可选的外部 compact 语义增强，用于在命名信息不足或命名风格差异较大时提供更丰富的对齐线索。当启用 \texttt{--use-compact} 且提供 \texttt{COMPACT\_EMB} 时，系统使用方法与测试外部 embedding 的余弦相似度替代名称重叠比例，并同样在所有测试上取最大值以保持与默认策略一致的聚合口径。该最大余弦相似度随后被映射到 $1\!\sim\!101$ 区间，使外部语义与默认语义在输出形式与尺度控制上保持一致，便于在相同模型结构与超参数设置下直接替换。若提供 \texttt{COMPACT\_MAP}，则使用 compact token 的重叠比例替代默认重叠，并沿用最大值聚合与整数映射的处理流程，从而保证不同语义来源在计算逻辑上具备可对照性。外部资源缺失或计算失败时，流程自动回退到默认名称重叠策略，以避免引入额外失败模式，并保证语义分数始终可获得。

\subsection{语义分数的落位方式与实现参数}

语义分数的落位方式遵循最小侵入原则，核心意图是让语义线索进入模型但不改变主干传播逻辑。实现中，语义分数不被设计为传播阶段的独立权重，也不引入语义专用消息通道，而是在节点初始化阶段作为 1 维特征直接拼接到节点嵌入中。方法节点与测试节点共享同一注入口径，它们的基础嵌入由 \texttt{token\_embedding} 给出，维度采用 \texttt{embedding\_size-1} 的设置。在当前配置下，\texttt{embedding\_size} 取 32，因此基础嵌入维度为 31，随后与 1 维语义分数 \texttt{inputtext} 拼接，得到 32 维的节点初始表示。这样的构造方式使语义信号以一维附加属性的形式进入节点表示空间，既显式可解释，又不会迫使模型为语义单独学习一套传播与融合规则。

代码行节点承担细粒度证据承载与历史线索补充的角色，其构造方式与方法测试节点不同。实现中，代码行节点使用 \texttt{token\_embedding1} 作为基础嵌入，其维度采用 \texttt{embedding\_size-2} 的设置。在 \texttt{embedding\_size}=32 的情况下，基础嵌入维度为 30，随后再拼接 \texttt{modification} 与 \texttt{churn} 两个标量特征，其中 \texttt{churn} 会被归一化到 $[0,1]$ 以抑制尺度漂移，而 \texttt{modification} 作为直接强度提示保留原值。这样一来，行节点携带的历史风险线索能够通过归属关系汇聚到方法节点，从而与方法测试的语义分数在方法候选空间内形成互补。

由于语义分数仅在初始化阶段注入，其在图中的扩散并非由语义专用机制决定，而是被后续结构学习模块间接塑造。结构学习通过对候选边的软选择抑制噪声边并保留关键通道，从而影响信息在图中的传播路径与强度。语义作为节点属性进入传播后，会随着节点表示一起被邻域聚合，但哪些邻域能有效参与聚合、哪些关系边应当被削弱或保留，最终由门控机制决定。在当前配置中，门控正则强度通过 \texttt{edge\_gate\_l1} 控制，默认设为 $1\times10^{-4}$，并通过 \texttt{edge\_gate\_warmup}=5 的渐进策略避免过早形成过强稀疏压力。门控输出还通过下界与平滑强度进行约束，\texttt{gate\_min}=0.6 保留基础结构通道，\texttt{gate\_strength}=0.2 控制门控对传播的影响幅度，温度参数 \texttt{gate\_tau}=0.5 用于调节门控分布的软硬程度，\texttt{gate\_bias\_init}=-2 则为门控初始状态提供更保守的起点。将这些设置融入传播机制后，语义信号既能在候选同质化时提供直接区分线索，又能在结构噪声存在时通过门控的筛边作用避免被无关通道稀释或错误扩散，从而与结构约束传播形成一致的工作方式。

综上，本文以名称相似度为主、外部 compact 语义为可选增强，构建了一条轻量语义线索。它从标识符规范化与分词开始，通过最大相似度聚合得到单一语义分数，再以一维特征的形式落位到节点初始表示中，并在后续传播中由结构学习门控间接决定其扩散路径与有效作用范围。该设计不追求对语义的复杂建模，而强调在候选打平时提供稳定的增量区分能力，并与现有结构传播主干保持低耦合的一致性。
